\documentclass[11pt]{article}
\usepackage[a4paper,margin=1in]{geometry}
\usepackage{amsmath,amssymb}
\title{Sharp Summary: Pauline et al. (2025) -- Foundations in General State Spaces}
\author{}
\date{}
\begin{document}
\maketitle

\section*{Scope}
A unifying tutorial/framework paper: discrete-time and continuous-time diffusion, continuous and discrete state spaces, and generator-based formalism in one notation.

\section*{Discrete diffusion equations (explicit)}
For categorical diffusion (single dimension), forward transition matrix $\tilde Q_t$ is
\[
q(x_t\mid x_{t-1})=\mathrm{Cat}(x_t;\,\tilde Q_t e_{x_{t-1}}),
\qquad
Q_{t\mid s}=\prod_{i=s+1}^t \tilde Q_i.
\]
Interpolation form:
\[
\tilde Q_t=\tilde\alpha_t I+(1-\tilde\alpha_t)p_{\text{noise}}\mathbf 1^\top,
\]
which yields
\[
q(x_t\mid x_0)=\mathrm{Cat}\big(x_t;\,\alpha_t e_{x_0}+(1-\alpha_t)p_{\text{noise}}\big),
\quad
\alpha_t=\prod_{j=1}^t\tilde\alpha_j.
\]

Reverse conditional given clean state:
\[
q(x_{t-1}\mid x_t,x_0)
=\mathrm{Cat}\!\left(
\frac{Q_{t\mid t-1}^\top e_{x_t}\odot Q_{t-1}e_{x_0}}
{e_{x_t}^\top Q_t e_{x_0}}
\right).
\]
This is exactly the algebra used in practical discrete reverse parameterizations.

\section*{Parameterization and training viewpoint}
The paper organizes training as approximation of
\[
q(x_{t-1}\mid x_t)
\approx p_\theta(x_{t-1}\mid x_t,t),
\]
with categorical/Gaussian families depending on state space, and derives ELBO-style objectives in both discrete and continuous time.

\section*{Key insight}
Its strongest value is structural: it cleanly separates
\begin{itemize}
\item forward design (kernel/rate choice),
\item reverse model class,
\item likelihood lower-bound optimization,
\end{itemize}
and shows how the same logic maps to CTMC and SDE settings.

\section*{Relevance to this repo}
\textbf{Direct and very high.}
Best source to keep notation consistent across notebook + tutorial text, especially for: forward kernels, Bayes reverse conditionals, and the interpretation of step-3 learning as reverse-conditional estimation across time.

\end{document}
